\begin{examplebox}
	\textbf{\textcolor{CExample}{例 1（基础）}}

	\textbf{\textcolor{CExample}{题目：}}
	已知数列 $\{a_n\}$ 满足 $a_{n+1}=2a_n+3$，$a_1=1$，求通项公式。

	\textbf{\textcolor{CExample}{解题步骤：}}
	\begin{description}[style=nextline, leftmargin=2.8em, labelsep=0.5em, itemsep=0.45em]
		\item[\textbf{第一步（求不动点）：}] 解方程 $\lambda=2\lambda+3$，得 $\lambda=-3$。
			\par\textbf{理由：}
			构造等比需要确定不动点。
		\item[\textbf{第二步（构造等比）：}] 令 $b_n = a_n - \lambda$，即 $b_n = a_n + 3$，则 $b_{n+1} = 2b_n$。首项 $b_1 = a_1 + 3$，得 $b_1 = 4$。
			\par\textbf{理由：}
			代入递推消去常数。
		\item[\textbf{第三步（求通项）：}] $b_n = 4 \cdot 2^{n-1}$，即 $b_n = 2^{n+1}$。所以 $a_n = b_n - 3$，故 $a_n = 2^{n+1} - 3$。
			\par\textbf{理由：}
			等比通项再减回 $\lambda$。
	\end{description}

	\textbf{\textcolor{CExample}{关键结论：}}
	\textbf{$a_n=2^{n+1}-3$}

	\medskip
	\textbf{\textcolor{CExample}{例 2（稍变形）}}

	\textbf{\textcolor{CExample}{题目：}}
	已知数列 $\{c_n\}$ 满足 $c_{n+1}=-c_n+4$，$c_1=1$，求通项公式。

	\textbf{\textcolor{CExample}{解题步骤：}}
	\begin{description}[style=nextline, leftmargin=2.8em, labelsep=0.5em, itemsep=0.45em]
		\item[\textbf{第一步（求不动点）：}] 解方程 $x=-x+4$，得 $\lambda=2$。
			\par\textbf{理由：}
			公式 $\lambda=\dfrac{q}{1-p}$ 或直接解。
		\item[\textbf{第二步（构造等比）：}] 令 $d_n = c_n - \lambda$，即 $d_n = c_n - 2$，则 $d_{n+1} = -d_n$。首项 $d_1 = c_1 - 2$，得 $d_1 = -1$。
			\par\textbf{理由：}
			消去常数。
		\item[\textbf{第三步（求通项）：}] $d_n = -1 \cdot (-1)^{n-1}$，计算得 $d_n = (-1)^n$。所以 $c_n = d_n + 2$，即 $c_n = 2 + (-1)^n$。
			\par\textbf{理由：}
			等比通项（公比 $p=-1$）加回 $\lambda$。
	\end{description}

	\textbf{\textcolor{CExample}{关键结论：}}
	\textbf{$c_n=2+(-1)^{n}$}

	\medskip
	\textbf{\textcolor{CExample}{例 3（综合应用）}}

	\textbf{\textcolor{CExample}{题目：}}
	已知数列 $\{a_n\}$ 满足 $a_{n+1}=2a_n+3$，$a_1=1$（同上），求其前 $n$ 项和 $S_n$。

	\textbf{\textcolor{CExample}{解题步骤：}}
	\begin{description}[style=nextline, leftmargin=2.8em, labelsep=0.5em, itemsep=0.45em]
		\item[\textbf{第一步（求通项）：}] 由例1得 $a_n=2^{n+1}-3$。
			\par\textbf{理由：}
			先求通项才能求和。
		\item[\textbf{第二步（拆分求和）：}] $S_n = \sum_{k=1}^{n}(2^{k+1}-3) = \sum_{k=1}^{n}2^{k+1} - \sum_{k=1}^{n}3$。
			\par\textbf{理由：}
			利用和的线性性质。
		\item[\textbf{第三步（计算求和）：}] $\sum_{k=1}^{n}2^{k+1}=2^{2}+2^{3}+\dots+2^{n+1}$。由等比求和公式得 $\sum_{k=1}^{n}2^{k+1}=2^{n+2}-4$。又 $\sum_{k=1}^{n}3=3n$，故 $S_n = 2^{n+2} - 3n - 4$。
			\par\textbf{理由：}
			等比数列求和公式。
	\end{description}

	\textbf{\textcolor{CExample}{关键结论：}}
	\textbf{$S_n=2^{n+2}-3n-4$}
\end{examplebox}
